\begin{theorem}[Elementary identities for Delta]\label{mod:delta-elementary}
Let $F$ be a nonarchimedean local field.  All local constants below are the
finite constants of Definition~\ref{mod:delta-finitedefinition}; admissible
denominators remain explicit, and denominator independence permits arbitrary
admissible choices in every displayed identity.

Let $\chi$ have exact multiplicative conductor $m$, let $\psi$ have exact
additive conductor $n$, and let $a\in F^\times$.  Define
$\psi^a(x)=\psi(ax)$.  Its exact conductor is
\[
 n_F(\psi^a)=n_F(\psi)+\operatorname{ord}_F(a).
\]
If $\gamma$ is admissible for $(\chi,\psi)$, then $a\gamma$ is admissible
for $(\chi,\psi^a)$.  At every numerator representative $u\in U_F^0$ one
has the literal identity
\[
 \psi^a\!\left(\frac{u}{a\gamma}\right)\chi(u)^{-1}
 =\psi\!\left(\frac{u}{\gamma}\right)\chi(u)^{-1}.
\]
After descending both sides to $U_F^0/U_F^m$ and summing, this gives
\[
 G_{a\gamma}(\chi,\psi^a)=G_\gamma(\chi,\psi),
 \qquad
 \Delta_F^{\mathrm{fin}}(\chi,\psi^a;\gamma_a)
 =\chi(a)\Delta_F^{\mathrm{fin}}(\chi,\psi;\gamma)
\]
for every admissible $\gamma_a$ for $(\chi,\psi^a)$.

For residual normalization, let $\bar\psi$ and $\psi_{\rm can}$ be
nontrivial complex-valued additive characters of $k_F$.  There is a local
unit $a\in\mathcal O_F^\times$ such that
\[
 \bar\psi\bigl(\overline a\,x\bigr)=\psi_{\rm can}(x)
 \qquad(x\in k_F).
\]
Formally, nontriviality makes
$c\mapsto\bar\psi(cx)$ a bijection from $k_F$ to its additive character
group; the scalar carrying $\bar\psi$ to $\psi_{\rm can}$ is nonzero and is
then lifted along the surjection
$\mathcal O_F^\times\twoheadrightarrow k_F^\times$.

Let $\mathbf 1$ denote the globally trivial quasi-character, packaged with
its exact conductor zero.  The quotient $U_F^0/U_F^0$ is a singleton.  For
an admissible $\gamma$ the sole finite summand is one, because
$1/\gamma\in\mathfrak p_F^{-n_F(\psi)}$.  Consequently
\[
 G_\gamma(\mathbf 1,\psi)=1,
 \qquad
 \Delta_F^{\mathrm{fin}}(\mathbf 1,\psi;\gamma)=1.
\]

Let $\mu$ be a finite-order quasi-character, and package $\mu^{-1}$ with
the same exact conductor.  For arbitrary admissible denominators $\gamma$
and $\gamma_{\rm inv}$,
\[
 \Delta_F^{\mathrm{fin}}(\mu,\psi;\gamma)
 \Delta_F^{\mathrm{fin}}(\mu^{-1},\psi;\gamma_{\rm inv})
 =\mu(-1).
\]
At a common denominator the representative-free finite sums satisfy the
exact conjugation formula
\[
 G_\gamma(\mu^{-1},\psi)
 =\mu(-1)\,\overline{G_\gamma(\mu,\psi)}.
\]
This is proved by the finite reindexing $u\mapsto-u$.  The conjugation step
uses two explicit unitary facts: quasi-character values on $U_F^0$ factor
through a finite quotient, and every additive-character value descends to
a finite lattice quotient and hence has norm one.  The two denominator
values cancel, while phase and conjugation give the displayed factor
$\mu(-1)$ with no additional sign.

Finally, let $S$ be a finite subgroup of the continuous quasi-characters of
$F^\times$, with $|S|$ an odd prime.  For each $\mu\in S$, choose exact
conductor data realizing its underlying character and an admissible
denominator $\gamma_\mu$.  Then
\[
 \prod_{\mu\in S}
 \Delta_F^{\mathrm{fin}}(\mu,\psi;\gamma_\mu)=1.
\]
Indeed every $\mu\in S$ has finite order dividing $|S|$, so $\mu(-1)=1$;
inversion pairs the nontrivial elements of $S$, and oddness makes the
identity its only fixed point.  The trivial factor is one by the preceding
formula.
\LeanFile{LanglandsFirstMainLemma/Delta/Elementary.lean}
\lean{LanglandsFirstMainLemma.delta_additive_scale}
\lean{LanglandsFirstMainLemma.residual_addChar_normalization}
\lean{LanglandsFirstMainLemma.delta_trivial_character}
\lean{LanglandsFirstMainLemma.delta_inverse_pair}
\lean{LanglandsFirstMainLemma.delta_odd_prime_character_product}
\leanok
\SourceText{Lemmas lem:additive-scale, lem:residual-normalization, lem:delta-one, lem:inverse-pair; Corollaries cor:additive-independence and cor:S-product-one.}
\uses{mod:delta-finitedefinition}
\FormalizationContract Additive scaling is proved before summation on the
same quotient representatives, and residual normalization returns an actual
local-unit lift of the nonzero residual scalar.  The inverse formula uses the
finite-order hypothesis stated in the manuscript and proves the required
unitarity and finite-sum conjugation internally.  The odd-prime product is
formulated for a finite subgroup of local quasi-characters so that the later
group $S(K/F)$ is an immediate instance; conductor packages and admissible
denominators are explicit dependent inputs.
\end{theorem}
